% eksperymentalne tłumaczenie na włoski

%

\section{Konstrukcje bez linijki} % lang-pl

\section{Costruzioni senza riga} % lang-it

% https://www.deltami.edu.pl/2015/08/prosto-w-srodek/ % TODO

Coxeter \cite[s. 95]{coxeter_1967} sugeruje skonstruować (samym) cyrklem wierzchołki sześciokąta foremnego, odcinek dwa razy dłuższy (i krótszy) od danego, obraz inwersyjny danego punktu leżącego dowolnie blisko środka i wreszcie podzielić dany odcinek na $n$ równych części. % lang-pl
Coxeter \cite[p. 95]{coxeter_1967} suggerisce di costruire con il solo compasso i vertici di un esagono regolare, un segmento lungo il doppio (e la metà) di un segmento dato, l'immagine inversiva di un punto dato situato arbitrariamente vicino al centro e infine di dividere un segmento dato in $n$ parti uguali. % lang-it


Okazuje się, że do żadnej klasycznej konstrukcji nie jest potrzebna linijka. % lang-pl
Si scopre che nessuna costruzione classica richiede realmente l'uso della riga. % lang-it

\begin{theorem}[twierdzenie Mohra-Mascheroniego] % lang-pl
\begin{theorem}[teorema di Mohr-Mascheroni] % lang-it
\index{twierdzenie!Mohra-Mascheroniego}%
    Jeśli dana konstrukcja jest wykonalna za pomocą cyrkla i~linijki, to jest ona wykonalna za pomocą samego cyrkla, o ile pominiemy rysowanie linii i ograniczymy się do wyznaczania punktów konstrukcji. % lang-pl
    
    Se una costruzione è realizzabile con riga e compasso, allora è realizzabile con il solo compasso, purché si ometta il tracciamento delle linee e ci si limiti a determinare i punti della costruzione. % lang-it
\end{theorem}
% PRZECZYTANO: https://en.wikipedia.org/wiki/Mohr-Mascheroni_theorem

Wynik będzie mieć ciekawą historię. % lang-pl

Pierwszy znajdzie go Georg Mohr, włoski geometra i~poeta w \emph{Euclides Danicus} (1672), ale jego odkrycie popadnie w zapomnienie, po czym będzie tak marnieć aż do roku 1928! % lang-pl
Questo risultato avrà una storia interessante. % lang-it

A scoprirlo per primo sarà Georg Mohr, geometra e poeta danese, nell'opera \emph{Euclides Danicus} (1672), ma la sua scoperta cadrà nell'oblio e vi rimarrà fino al 1928! % lang-it
% established his results by using the idea of reflection in a line.
% Euclides Danicus = duński Euklides
\index[persons]{Mohr, Georg}%
Niezależnie twierdzenie odkryje Lorenzo Mascheroni; wykorzysta odbicia względem prostych i~opisze je w~\emph{La Geometria del Compasso} (1797). % lang-pl

Indipendentemente, il teorema sarà riscoperto da Lorenzo Mascheroni, che utilizzerà le riflessioni rispetto alle rette e le descriverà nella sua opera \emph{La Geometria del Compasso} (1797). % lang-it
\index[persons]{Mascheroni, Lorenzo}%
W 1890 jeszcze raz to samo, ale nie tak samo (inwersjami) zrobi śląski\footnote{Urodzony w Księstwie Górnego i Dolnego Śląska, znanym także jako Śląsk Austriacki} matematyk August Adler w pracy, której tytuł nadgryzie ząb czasu. % lang-pl

Nel 1890 il matematico slesiano\footnote{Nato nel Ducato dell'Alta e Bassa Slesia, noto anche come Slesia austriaca.} August Adler otterrà nuovamente lo stesso risultato, ma con un metodo diverso, basato sulle inversioni, in un lavoro il cui titolo è stato ormai consumato dal tempo. % lang-it
% August Adler (24 January 1863, Opava, Austrian Silesia – 17 October 1923, Vienna) was a Czech and Austrian mathematician noted for using the theory of inversion to provide an alternate proof of Mascheroni's compass and straightedge construction theorem. Czyli nie był z niczego innego znany?
Wreszcie w 1928 student Hjelmsleva przeglądając księgarnię w Kopenhadze i jej zawartość znajdzie książkę Mohra. % lang-pl

Infine, nel 1928, uno studente di Hjelmslev troverà il libro di Mohr sfogliando i volumi di una libreria di Copenaghen. % lang-it
% źródło tej rewelacji: George Martin - Geometric Constructions (1998), s. 54
Będzie to duża niespodzianka dla Hjelmsleva. % lang-pl
Sarà una grande sorpresa per Hjelmslev. % lang-it


Dla dowodu wystarczy pokazać, że cyrkel wystarcza do przeprowadzenia pięciu konstrukcji: % lang-pl
Per la dimostrazione basta mostrare che il compasso è sufficiente per eseguire le seguenti cinque costruzioni: % lang-it
\begin{enumerate}
    \item poprowadzenia prostej przez dwa punkty, % lang-pl
    \item wykreślenia okręgu o danym środku i promieniu, % lang-pl
    \item znalezienia punktu przecięcia dwóch nierównoległych prostych, % lang-pl
    \item znalezienia punktu (lub punktów) przecięcia prostej z okręgiem, % lang-pl
    \item znalezienia punktu (lub punktów) przecięcia dwóch okręgów. % lang-pl
    \item tracciare la retta passante per due punti, % lang-it
    \item tracciare un cerchio di centro e raggio assegnati, % lang-it
    \item trovare il punto di intersezione di due rette non parallele, % lang-it
    \item trovare il punto (o i punti) di intersezione tra una retta e un cerchio, % lang-it
    \item trovare il punto (o i punti) di intersezione di due cerchi. % lang-it
\end{enumerate}

O wyniku Mohra-Mascheroniego napisze Eves \cite[s. 169-172]{eves1_1972}. % lang-pl

Del risultato di Mohr-Mascheroni scrive Eves \cite[p. 169-172]{eves1_1972}. % lang-it

%

\begin{problem}[zadanie Napoleona] % lang-pl
\begin{problem}[problema di Napoleone] % lang-it
\index{zadanie!Napoleona}
\label{napoleon_problem}
	Dany jest okrąg oraz jego środek. % lang-pl
	
	Podzielić okrąg na cztery łuki równej miary korzystając z cyrkla, ale nie linijki. % lang-pl
	Sono dati un cerchio e il suo centro. % lang-it
	
	Dividere il cerchio in quattro archi di uguale ampiezza usando il compasso ma non la riga. % lang-it
\end{problem}
% PRZECZYTANO: https://en.wikipedia.org/wiki/Napoleon%27s_problem

Jak zawsze w przypadku Napoleona i geometrii, nie wiadomo, czy rozwiązał albo choć wymyślił przedstawione wyżej zadanie konstrukcyjne. % lang-pl

Come spesso accade quando si parla di Napoleone e geometria, non si sa se abbia risolto o persino formulato il problema di costruzione sopra descritto. % lang-it
\index[persons]{Bonaparte, Napoleon}%
W trudniejszej wersji okrąg nie ma środka i trzeba go najpierw znaleźć. % lang-pl

Rozwiązanie podają Bogdańska, Neugebauer \cite[s. 116]{neugebauer_2018} (z wykorzystaniem okręgów Torricelliego); % lang-pl
Nella versione più difficile il cerchio non ha il centro indicato e bisogna trovarlo per prima cosa. % lang-it

u Audina \cite[s. 105]{audin_2003} jest to ćwiczenie. % lang-pl
Una soluzione è fornita da Bogdańska e Neugebauer \cite[p. 116]{neugebauer_2018} (utilizzando i cerchi di Torricelli); % lang-it

in Audin \cite[p. 105]{audin_2003} compare invece come esercizio. % lang-it
\index{okrąg Torricelliego}%

%
